\chapter*{Reform and Revolution}
\addcontentsline{toc}{chapter}{Reform and Revolution}
\setcounter{postnote}{0} % Restarting endnotes numbering for each chapter

\section*{1}
At Marx's graveside on 17 March 1883, Engels said of him that he was 'before all else a revolutionist' and that 'his real mission in life was to contribute, in one way or another, to the overthrow of capitalist society and of the state institutions which it had brought into being, to contribute to the liberation of the modern proletariat\dots'\postnote{SW 1968, p. 436.} Commonplace though this may be, it is very much worth stressing that this was indeed Marx's fundamental and unwavering purpose; and that Marxism as a political doctrine has above all been about the making of socialist revolution.

But what strategy does the making of socialist revolution require? A hundred years of debate have shown well enough that however unambiguous the Marxist purpose may be, its advancement has been by far the most contentious of all issues within the ranks of Marxism. Nor has the debate grown less intense with the years and with the accumulated experience they have provided. The lines of division have remained as deep as ever on this central issue of strategy, and will undoubtedly remain so for a long time to come.

The reason for the endurance and intractability of these divisions is that they represent two very different paths of advance, and the differences also concern the meaning which is to be attached to the notion of socialist 'revolution' and the 'overthrow' of capitalist regimes. Each of these different paths have appeared to their respective champions as the only reasonable and realistic one, by which token the alternative one has been condemned as unrealistic, self-defeating and a betrayal or deformation of Marxism; and it has also been possible for the respective antagonists to defend their chosen positions, with varying degrees of plausibility, as representing authentic Marxist positions, or as representing the only possible Marxist position in the circumstances, or some such argument.

It is not pointless or impossible to determine fairly precisely what Marx's position on this issue of strategy was at any given time, or Engels's, or Lenin's. But this never settles anything; and  the more important requirement, in the present context, is to try and clarify what the real lines of division are on the subject of revolutionary strategy within the Marxist perspective. It is particularly necessary to try and do this because the inevitable acerbity and virulence of the controversies have greatly obscured what these lines of division actually are; and also what they are not.

The two paths in question are in essence those which have commonly been designated as the 'reformist' and the 'revolutionary'. For reasons which will be discussed presently, these terms are rather misleading and fail to provide a proper contraposition of the alternatives. Those Marxists who have been called 'reformists' by their 'revolutionary' Marxist opponents have vehemently repudiated the label as an insult to their own revolutionary integrity and purpose; and they have in turn denounced their opponents as 'adventurist', 'ultra-left', etc. All the designations which have formed part of this Marxist debate on strategies of advance have been greatly overloaded with question-begging connotations, and no label in this controversy is altogether 'innocent'. 'Constitutionalist' and 'insurrectionary' probably come nearest to an accurate description of the alternative positions; but there are problems here too. I have continued to use the term 'reformist', in a sense that will be specified, and without pejorative intent.

Before proceeding with the discussion of the alternative strategies, there is one source of confusion which must be cleared up. This is the fact that there has always existed a trend in working-class movements—and for that matter outside—towards social reform; and this is a trend which, in so far as it has no thought of achieving the wholesale transformation of capitalist society into an entirely different social order, must be sharply distinguished from the 'reformist' strategy, which has insisted that this was precisely its purpose.

As was noted earlier, social reform has been an intrinsic part of the politics of capitalism, and those who have supported it have not only not been concerned to advance towards socialism but have on the contrary seen in social reform an essential prophylactic against it. It was this which Marx and Engels called in the Communist Manifesto 'Conservative or Bourgeois Socialism' and which they attributed to 'a part of the bourgeoisie\dots desirous of redressing social grievances, in order to secure the continued existence of bourgeois society'.\postnote{Revs., p. 93.} 'The socialistic bourgeois', they also wrote, 'want all the advantages of modern social conditions without the struggles and dangers necessarily resulting therefrom. They desire the existing state of society minus its revolutionary and disintegrating elements. They wish for a bourgeoisie without a proletariat'; and the reforms that such people wanted by no means involved the 'abolition of the bourgeois relations of production', which required a revolution, and in no respect affect the relations between capital and labour'.\postnote{Ibid., pp. 92-3.}

The last statement is obviously an exaggeration, since reforms can affect and have affected in some respects the relations between capital and labour; and both Marx and Engels worked very hard to achieve reforms to that end, for instance inside the First International.\postnote{See e.g. H. Collins and C. Abramsky, Karl Marx and the British Labour Movement (London, 1965).} But the characterization of 'Conservative or Bourgeois SocialismTemains generally speaking extraordinarily modern, notwithstanding the passage of 130 years; and it encompasses much more than 'part of the bourgeoisie'. Marx and Engels themselves saw the same limited and prophylactic intentions in the reforms of straightforward conservative figures such as Louis Bonaparte (whose 'Imperial Socialism' Marx mocked), Bismarck and Disraeli; and they also applied it, in the aftermath of the 1848 revolutions, to what they called the 'republican petty bourgeois', 'who now call themselves Red and social-democratic because they cherish the pious wish of abolishing the pressure of big capital on small capital, of the big bourgeois on the small bourgeois'. Such people, they said, 'far from desiring to revolutionize all society for the revolutionary proletarians\dots strive for a change in social conditions by means of which existing society will be made as tolerable and comfortable as possible for them'.\postnote{K. Marx and F. Engels, 'Address of the Central Committee to the Communist League', in SW 1950, I, p. 101.}

But Marx and Engels also saw develop in their lifetime movements and associations of workers whose aims did not go beyond the achievement of specific and limited reforms: British trade unionism is a case in point, and Marx and Engels frequently deplored its narrow and firmly un-revolutionary consciousness and purposes,\postnote{See e.g. K. Marx and F. Engels, On Britain (Moscow, 1953), passim.} which did not however stop them from working with British trade unionists when they found this useful to their own purposes, as in the First International.\postnote{See e.g. Collins and Abramsky, op. cit.} It is also this concern with limited amelioration which Lenin described as 'trade union consciousness', and which he wanted  the revolutionary party to extend to larger anti-capitalist purposes. Of course trade unions, however limited their horizons, have seldom been able to concern themselves exclusively with hours, wages and conditions' if only because these tend to bring other issues into focus. But this need not go beyond improvements of various kinds for the working class within the capitalist system, and with no thought of superseding that system; and this trade union consciousness' has been and remains an important and in some instances a dominant trend inside organized labour under capitalism.

In the twentieth century this trend has not only been present in the trade unions: it has also come to dominate large (and in some cases the largest) working-class parties in capitalist countries. Britain, Germany, and Sweden are obvious instances. No doubt most such parties have a formal and explicit commitment to the achievement of a different social order; and they also include many people who deeply believe in the reality and meaningfulness of that commitment on the part of their parties—or who believe at least that their parties can eventually be made to take that commitment seriously.

Whatever may be thought of that hope, it must be obvious that the parties concerned are in fact parties of social reform, whose leaders and leading personnel are in their overwhelming majority solidly and comfortably established in the existing social order, and who have absolutely no intention of embarking on anything resembling its wholesale transformation, in however piecemeal and pacific a perspective. Their purpose is reform, very often conceived, in more or less the same way as it is by their conservative counterparts, as a necessary insurance against pressure for reforms of a too radical and rapid kind. The 'socialism' which these leaders proclaim, where they do proclaim it, is a rhetorical device and a synonym for various improvements that a necessarily imperfect society requires. These reforms do not form part of a coherent and comprehensive strategy of change, least of all socialist change.

'Reformism', on the contrary, is such a strategy, at least theoretically. Indeed, not only is it one of the two main strategies of the Marxist tradition: for good or ill, it is also the one which, notwithstanding a deceptive rhetoric, has always found most favour and support within that tradition. This strategy naturally includes the pursuit of reforms of every kind - economic, social, and political—within the framework of capitalism. But unlike parties of social reform, parties which are guided by this particular Marxist tradition do not consider such reforms as being their ultimate purpose: these are at best steps and partial means towards a much larger purpose, which is declared to be the 'overthrow' of capitalism and the achievement of an altogether different, that is socialist, society.

I have already noted in Chapter IV how easily Marx and Engels accepted the idea that bourgeois democratic regimes afforded the most promising ground for the development of the proletarian revolutionary movement; and also how easy they found it to envisage a strategy for the revolutionary movement which entailed thorough involvement in 'ordinary' political life and the pursuit of reform, allied to a constant concern for the advancement of a revolutionary purpose that naturally went far beyond anything that reformers of one sort or another could envisage.

The distinctiveness of that Marxist revolutionary purpose was emphasized again and again by Marx and Engels. In The Class Struggles in France, Marx thus identifies it as
\begin{displayquote}
    the declaration of the permanence of the revolution, the class dictatorship of the proletariat as a necessary intermediate point on the path towards the abolition of class differences in general, the abolition of all social relations which correspond to these relations of production, and the revolutionizing of all ideas which stem from these social relations.\postnote{SE, p. 123.}
\end{displayquote}

Similarly in the Address to the Central Committee to the Communist League to which reference was made earlier, Marx and Engels insisted that the limited demands of the 'democratic petty bourgeois' could in no way suffice 'for the party of the proletariat':
\begin{displayquote}
    It is our interest and our task to make the revolution permanent, until all more or less possessing classes have been forced out of their position of dominance, until the proletariat has conquered state power, and the association of proletarians, not only in one country but in all the dominant countries of the world, has advanced so far that competition among the proletarians of these countries has ceased and that at least the decisive productive forces are concentrated in the hands of the proletarians. For us the issue cannot be the alteration of private property but only its annihilation, not the smoothing over of class antagonisms but the abolition of classes, not the improvement of existing society but the foundation of a new one.\postnote{Address ... in SW 1950, I, p. 102.}
\end{displayquote}

Making the revolution permanent', in this context, clearly  means striving for the advancement of these aims within the framework of capitalism and a bourgeois democratic regime; and this striving obviously included pressure for reforms of every sort. Indeed, Lenin in 1905 made the achievement of a bourgeois democratic republic the immediate aim of the revolutionary movement. In Two Tactics of Social Democracy in the Democratic Revolution, he said that 'we cannot get out of the bourgeois-democratic boundaries of the Russian revolution, but we can vastly extend these boundaries, and within these boundaries we can and must fight for the interests of the proletariat, for its immediate needs and for conditions that will make it possible to prepare its forces for the future complete victory.'\postnote{Two Tactics \dots{} in SWL, p. 78.}

This perspective entails a combination of two concepts which have often been erroneously contraposed in Marxist thinking on strategies of revolutionary advance, namely the concept of 'two stages' and that of 'permanent revolution'. The 'two stages' concept means in effect the struggle for the achievement of a bourgeois democratic regime where such a regime does not exist, and where conditions for a socialist revolution are deemed not to exist either. The concept of 'permanent revolution', at least in Marx's usage, involves the application of continuous pressure for pushing the process further and creating the conditions for a revolutionary break with the bourgeois democratic regime itself.\footnote{As envisaged by Trotsky, with whom the notion of 'permanent revolution' is most closely associated, it entailed a much more dramatic belief in the need and possibility to by-pass in Russia the bourgeois revolution and move directly to the dictatorship of the proletariat; and 'permanent revolution' also meant for him the fostering of revolution abroad as an essential part of the process. (See Deutscher, The Prophet Armed, op. cit., Ch. VI.)} In 1905, Lenin was arguing not only that Marxists should struggle for the achievement of such a regime, but that they must assume the leadership of that struggle by way of the 'revolutionary-democratic dictatorship of the proletariat and the peasantry'.\postnote{Ibid., p. 82.} For all its vast significance, he also said, 'the democratic revolution will not immediately overstep the bounds of bourgeois social and economic relations.' But its success would 'mark the utter limit of the revolutionism of the bourgeoisie, and even that of the petty bourgeoisie, and the beginning of the proletariat's real struggle for socialism'.\postnote{Ibid., p. 82.}

Whatever may be thought of this 'scenario' (and the notion of a 'revolutionary\hyp{}democratic dictatorship of the proletariat and the peasantry' carefully refusing to overstep the bounds of bourgeois social and economic relations is certainly strange),\footnote{With the outbreak of revolution in Russia in 1917, Lenin abandoned this schema and his views as to what could and should be done moved much closer to those of Trotsky in 1905.} it is clearly the case that the struggle for reforms in a bourgeois democratic regime was never taken by classical Marxism to be incompatible with the advancement of revolutionary aims and purposes. On the contrary such a struggle is an intrinsic part of the Marxist tradition. With in that tradition, there is undoubtedly room for much controversy and debate as to the kind of reforms to be pursued, the importance that should be attached to them, and the manner in which they should be pursued. Thus supporters of a 'revolutionary' as opposed to a 'reformist' strategy have generally tended to place relatively less emphasis and value on reforms, and to press for reforms which they did not believe to be attainable, as part of a 'politics of exposure' of capitalism—and also of 'reformist' labour leaders. But there are limits to this kind of politics, which are imposed inter alia by the working\hyp{}class movement itself: if the 'politics of exposure' are pushed too far, all that they are likely to expose are the people who practice them, and leave such people in a state of ineffectual sectarian isolation. In any case, and notwithstanding the undoubted differences of emphasis within the Marxist tradition on the degree of importance to be attached to reforms, it is not this which essentially distinguishes the two contrasting strategies within that tradition. It is not, in other words, the pursuit of reforms which defines 'reformism'.

Nor is 'reformism', in its Marxist version, to be defined as gradualism', according to which the achievement of a socialist society is conceived as a slow but sure advance by way of a long sequence of reforms, at the end of which (or for that matter in the course of which) capitalism would be found to have been transcended. This is what Sidney Webb and the original Fabians roughly had in mind when they thought of socialism. But they were not Marxists, and loudly proclaimed that they were not; and their version of socialism had much more to do with piecemeal collectivist social engineering, inspired, directed, and administered from on high than with any version of the Marxism which they opposed. (It was not simply an aberration which led the Webbs to be so deeply attracted to the Soviet Union in the early thirties—there was much about it which matched their vision of socialism as state-imposed collectivism.) Whatever else may be said about 'reformist' Marxism, it never at any time envisaged the 'transition to socialism' in such a smoothly gradualist perspective, with so marked a stress as in the case of Fabianism on the conversion of members of the middle (and upper) classes to state intervention and collectivist measures.

Marxist 'reformism' does have a long-term view of the advance towards socialism, and that view does include a belief in the need to chip away at the structures of capitalism. But “reformism' envisages this process in terms of struggle and more specifically class struggle on many different fronts and at many different levels: in this sense, it remains quite definitely a politics of conflict.

The really important point, however, is that this is a politics of conflict envisaged as being conducted within the limits of constitutionalism defined by bourgeois democracy, with a strong emphasis on electoral success at municipal, regional and national levels, and with the hope of achieving majority or at least strong representation in local councils, regional assemblies, and national parliaments. Where relevant, this would also include competition in presidential elections, with the hope of either winning outright, or at least of being able to achieve a favourable bargaining position for this or that purpose: there are obviously many possible permutations.

This emphasis on constitutionalism, electoralism, and representation is certainly crucial in the definition of 'reformism'. But it is often caricatured by its opponents on the far left as a necessarily exclusive concern with electoral success and increased representation. In fact, 'reformism' is also compatible, both theoretically and practically, with forms of struggle which, though carried on within the given constitutional framework, are not related to elections and representation—for instance industrial struggles, strikes, sit-ins, work-ins, demonstrations, marches, campaigns, etc., designed to advance specific or general demands, oppose governmental policies, protest against given measures, and so on.\footnote{There is an interesting discussion of this strategy, and of its problems, in an article of 1904 by Rosa Luxemburg, entitled 'Social Democracy and Parliamentarism': Parliamentarism, she wrote, 'is for the rising working class one of the most powerful and indispensable means of carrying on the class struggle. To save bourgeois parliamentarism from the bourgeoisie and use it against the bourgeoisie is one of Social Democracy's most urgent political tasks'. However, she categorically rejected any abandonment of the extra-parliamentary class struggle: The real way is not to conceal and abandon the proletarian class struggle, but the very reverse: to emphasize strongly and develop this struggle both within and without parliament. This includes strengthening the extra-parliamentary action of the proletariat as well as a certain organization of the parliamentary action of our deputies' (Rosa Luxemburg. Selected Political Writings, R. Looker, Ed. (London, 1972), pp. 110, 113.) See also The Mass Strike, The Political Party and the Trade Unions, in Rosa Luxemburg Speaks, op. cit., p. 158.} No doubt, such activities and the manner or moment in which they are conducted may be influenced by electoral preoccupations, even of a fairly distant kind. But this is not always the case, and may not be the case at all.

This having been said, it is nevertheless right to stress the constitutionalism, electoral concerns, and representative ambitions of 'reformist' Marxist parties; and also to note the very large fact that it is parties of this kind and with these characteristics which have dominated working-class movements throughout the history of the bourgeois constitutional regimes of advanced capitalism—the main alternative being social reform parties like the British Labour Party. The reasons for this predominance are central to the nature of working-class politics and highlight some of the major problems which confront Marxist parties in this kind of political system.

The basic point is that bourgeois democracy and constitutionalism generate considerable constraints for revolutionary movements and lead them towards what might be termed reciprocal constitutionalism. Working-class movements in the shape of trade unions and parties which grew up in the shadow of capitalism and gained legal recognition and political acceptance could truthfully declare that their purpose was the achievement of a socialist society; and they could also proclaim their allegiance to Marxism. But they were nevertheless part of a functioning political system and their political mode of operation had to be legal and constitutional; or rather it was so, since Marxist parties do not, in such regimes and under 'normal' circumstances, choose illegality, un-constitutionalism, and clandestinity.\footnote{This needs some slight, qualification for the history of Western Communist parties in the first years of their existence, for which see below.}

Legality and constitutionality do not in themselves, at least in non-revolutionary circumstances, mean an abandonment of revolutionary purposes or need not necessarily do so. After all, there are in all capitalist countries with bourgeois democratic regimes parties and groupings which claim Marxist credentials and which advocate revolutionary policies, yet which operate within the framework of bourgeois legality. With greater or lesser enthusiasm, and with occasional confrontations with the forces of law and order, they work within the system: in this respect, the difference between 'reformist' parties and their opponents on the left is a matter of emphasis and perspective rather than of fundamental and immediate choice. Much more important is the difference created between them by the electoral ambitions of 'reformist' parties, and what this entails in terms of policies, programmes, and political behaviour.

Parties with serious electoral ambitions, however genuine their ultimate intention to form and transcend capitalist structures, are inevitably tempted to try and widen their appeal by emphasizing the relative moderation of their immediate (and not so immediate) aims. It was all very well for Engels, in his 1895 Introduction to Marx's Class Struggles in France which I have already quoted earlier, to say that the way things were going for German Social Democracy, 'we shall conquer the greater part of the middle strata of society, petty bourgeois and small peasants', as well as the bulk of the working-class vote.\postnote{SW 1950, I, p. 124; p. 665, in SW 1968.} One thing this ignored was that such a 'conquest' might have to be bought or achieved at a substantial cost to both ideology and political programme. Engels appeared to take it for granted on the contrary that the party's programme need not be diluted in order eventually to appeal to a majority of the German voters, working class and non-working class. But this was at best a very dubious assumption, which might come to be warranted but which party leaders would themselves regard as very dubious. This being the case, they would find all but irresistible the temptation to dilute their programmes in the direction of moderation and reassurance.

Secondly, Engels's perspective also ignored the temptation to which the leaders of mass working-class parties such as the German Social Democratic Party would be exposed, namely to see themselves as the custodians of vast organizations whose effectiveness and future prospects must not be put at risk by 'extreme' policies. He had warned against the frittering away of the 'shock force' constituted by the Party in what he called 'vanguard skirmishes', and insisted that their main task was to keep the party 'intact until the decisive day'.\postnote{Ibid.} But this could easily be translated (and was translated) into a considerable reluctance to 'fritter away' potential electoral support by policies and attitudes which might frighten voters off—and this of course meant blunting the edges of the Party's commitments. As for the 'decisive day', it was still sufficiently far off, and blurred enough in its meaning, to constitute no embarrassing point of reference to Party leaders.

On the basis of what has been said so far, Marxist 'reformism' entails first a concern with the day-to-day defence of working- class interests and the advancement of reforms of every kind; and secondly a thorough involvement in the politics of bourgeois democracy, with the intention of achieving the greatest possible degree of electoral support, and participation in parliamentary and other representative institutions—local council, regional assemblies, and the like.

If this was all that was meant by 'reformism', the debates and feuds which it has engendered in the ranks of Marxism would not have been nearly so bitter and irreconcilable. Undoubtedly, there would still have been much to argue about, and very sharply, but more in terms of strategy and tactics in relation to specific episodes, in terms of emphasis and attitudes, than of radically different over-all positions and general strategies. Indeed, this is precisely how matters stood in the great controversies between Marxists before 1914; and one of the features of Stalinism in the intellectual and historical fields was to read later events, positions, and attitudes backward in time, and to transpose the sharpness of the divisions of a later period into a previous one, when these divisions were of a much less accented character.

Thus much had already come to separate Lenin from Kautsky and other leaders of German Social-Democracy before 1914. But Kautsky, to take him as the most obvious example of the point, had nevertheless been emphatically opposed to Bernstein's 'revisionism' at the turn of the century; and even though his opposition was not nearly as fundamental as then appeared, Kautsky remained part of a European Marxist spectrum and was indeed one of the most eminent and respected figures of European Marxism, which of course included Russian Marxism. In  no way was he then in Lenin's eyes the 'renegade Kautsky' of later years. On the contrary, Lenin had in 1902 quoted Kautsky in What is to be Done? with much deference; and in 1905, he had rhetorically asked 'what and where did I ever claim to have created any sort of special trend in International Social- Democracy not identical with the trend of Bebel and Kautsky? When and where have there been brought to light differences between me, on the one hand, and Bebel and Kautsky, on the other . . .?' 'The complete unanimity of international revolutionary Social-Democracy on all major questions of programme and tactics', he had himself answered, 'is a most incontrovertible fact'. Lenin rather exaggerated 'the complete unanimity' of the international revolutionary movement even then; but the denial of fundamental differences between himself and his later enemies is nevertheless worth noting.\footnote{V. I. Lenin, Two Tactics\dots{}, in SWL, pp. 89-90. However, Lenin sharply reproved Kautsky in the same work for the latter's rather superior attitude to the quarrels of Russian revolutionaries on the position they shold adopt in regard to an eventual provisional revolutionary government. 'Although Kautsky, for instance', Lenin wrote, 'now tries to wax ironical and says that our dispute about a provisional revolutionary government is like sharing out the meat before the bear is killed, this irony only proves that even clever and revolutionary Social-Democrats are liable to put their foot in it when they talk about something they know only about by hearsay (ibid., p. 122). The point is sharply made, but there is none of the bitter enmity of later controversies. Nor was there reason to be—Lenin was criticising people whom he considered 'revolutionary Social-Democrats'.}

Leninism was not a revolutionary strategy hostile to parliamentary participation as such, either before 1914 or after. One of Lenin's most famous and influential pamphlets, 'Left-Wing' Communism—An Infantile Disorder, written in 1920, was in part directed against what he viewed as a sectarian and ultraleft deformation among Communists in Germany and other Western countries. In effect, Lenin was attacking the obverse of what Marx had called 'parliamentary cretinism', namely anti- parliamentary cretinism, which would prevent revolutionaries from making such usage for their own purposes as bourgeois parliaments did offer. Marx's own quarrels with the anarchists in the First International had had to do among other things with the contempt that many of them had for 'ordinary' politics and their rejection of sustained and disciplined organization by revolutionaries for the purpose of involvement in such politics.

Marx himself had no qualms whatever in being thus involved. Nor had Lenin.

The fundamental contraposition between 'reformism' and its alternative within Marxism, which may properly be called Leninism, concerned altogether different issues, of crucial and enduring importance.

For a start, the point needs to be made that Leninism after 1914 did become that 'sort of special trend' which Lenin, in 1905, denied having created in International Social Democracy. With the war, the collapse of the Second International and the endorsement by all of its principal constituents of the 'national interest', International Social Democracy as a more or less unified movement, or at least as one movement that could accommodate many great differences, was irreparably shattered. Lenin believed that the war had opened the era of intensified class conflict and placed proletarian revolution on the agenda; and that a new and very different organization from the previous one must bring together revolutionary parties of a very different sort from those that had dominated the Second International. In effect, what Lenin did as a result of the war was to place insurrectionary politics on the agenda, first for Russia, and then, with the immense prestige conferred upon him and the Bolsheviks by their conquest of power, for the international revolutionary movement.

\section*{2}
Insurrectionary politics' (the term is used here because 'revolutionary politics' is too loose) is not intended to suggest that Lenin wanted revolutionaries everywhere to prepare for immediate insurrection; and, as I have just noted, he did not believe that Communists could, even at this time, neglect politics of a different sort. It is true that he vastly exaggerated the revolutionary possibilities which existed in the aftermath of the war, and was no doubt greatly encouraged to do so by the revolutionary eruptions which had occurred in Germany, Hungary, and Austria, as well as by the very marked radical temper which had gripped large sections of the working class everywhere in 1918-20 and led to great industrial strikes and social agitation. But his misjudgements, which had very many negative consequences, nevertheless did not lead him to believe in instant revolution, made by will and proclamation. What he did  believe was that the time was ripe for preparing as a matter of extreme urgency for a seizure of power which might not immediately be possible in all capitalist countries but whose appearance on the revolutionary agenda could not be long delayed in many if not most of them.

What this meant is perhaps best illustrated by extensive quotation from the famous Twenty-One Conditions of admission to the Communist International, which were adopted at the Second World Congress of the International in Moscow in July\hyp{}August 1920.

The first such condition was that the daily propaganda and agitation of every party (now to be called 'Communist Party')
\begin{displayquote}
    must bear a truly communist character and correspond to the programme and all the decisions of the Third International. All the organs of the press that are in the hands of the party must be edited by reliable communists who have proved their loyalty to the cause of the proletarian revolution. The dictatorship of the proletariat should not be spoken of simply as a current hackneyed formula; it should be advocated in such a way that its necessity should be apparent to every rank-and-file working man and woman, each soldier and peasant, and should emanate from the facts of everyday life systematically recorded by our press day after day.
\end{displayquote}

The second condition required the removal from all 'responsible' posts 'in the labour movement, in the party organization, editorial boards, trade unions, parliamentary fractions, cooperative societies, municipalities, etc., all reformists and followers of the “center”, and have them replaced by Communists\dots{}'

The third condition declared that 'the class struggle in almost all countries of Europe and America\index{America} is entering the phase of civil war'. Under such conditions, it went on,
\begin{displayquote}
    the communists can have no confidence in bourgeois law. They must everywhere create a parallel illegal apparatus, which at the decisive moment could assist the party in performing its duty to the revolution. In all countries where, in consequence of martial law or exceptional laws, the communists are unable to perform their work legally, a combination of legal and illegal work is absolutely necessary.
\end{displayquote}

Subsequent conditions stipulated the duty of Communists to spread their propaganda in the countryside, to engage in unremitting struggle against reformists and 'half-reformists' everywhere, and to denounce imperialism and colonialism, this duty being particularly imperative for Communists in imperialist countries.

Condition 11 required Communist parties to
\begin{displayquote}
    overhaul the membership of their parliamentary fractions, eliminate a unreliable elements from them, to control these fractions, not only verbally but in reality, to subordinate them to the central committee of i and demand from every communist member of parliament that he devote his entire activities to the interests of really revolutionary propaganda and agitation.
\end{displayquote}

Condition 12 stipulated that parties belonging to the Third International 'must be built up on the principle of democratic centralism':
\begin{displayquote}
    Present time of acute civil war, the communist party will only be able fully to do its duty when it is organized in the most centralized manner, it it has iron discipline, bordering on military discipline, and if the party center is a powerful, authoritative organ, with wide powers possessing the general trust of the party membership.
\end{displayquote}

Condition 14 laid down that Communist parties 'must give every possible support to the Soviet Republics in their struggle against all counterrevolutionary forces', and that they should carry on propaganda to induce workers to refuse 'to transport munitions of war intended for enemies of the Soviet Republic'; and carry on 'legal and illegal propaganda among the troops which are sent to crush the workers' republic'.

The document made it absolutely clear that the new International would be organized on very different lines from its predecessors (and, it might be added, from the one before that—Marx's First International):
\begin{displayquote}
    All decisions of the congresses of the Communist International [Condition 16 stated] as weli as the decisions of its Executive Committee, are binding on all parties affiliated to the Communist International. The Communist International, operating in the midst of a most acute civil war, must have a far more centralized form of organisation than that of the Second International.
\end{displayquote}

As for the required change of name to 'Communist Party' (Condition 17), the document declared that this was not merely a formal question, but a political one of great importance':
\begin{displayquote}
    The Communist International has declared a decisive war against the entue bourgeois world and all the yellow, social democratic parties. Every rank-and-file worker must clearly understand the difference be-  tween the communist parties and the old official 'social democratic' or socialist parties which have betrayed the cause of the working class.
\end{displayquote}

The last of the Twenty-One Conditions perhaps best typifies the spirit of the whole document: 'Members of the party who reject the conditions and theses of the Communist International, on principle, must be expelled from the party.'\postnote{H. Gruber, International Communism in the Era of Lenin, pp. 241-6 for the full text.}

In so far as the Twenty-One Conditions made any sense at all, they did so as the battle orders of an international army, organized into a number of national units under a supreme command, and being prepared for an assault that could not be long delayed on the citadels of world capitalism; and that army was therefore having its ranks cleared of those tainted and untrustworthy elements who refused to purge their past mistakes and fall in line with the strategy presented to them by the leadership of the victorious Russian Revolution.

But the tide was already ebbing by 1920. Such international army as there had been was already beginning to disband by the time the Twenty-One Conditions were issued; and some of its main units had recently been decisively defeated. The new Communist parties, far from gaining a commanding position in their respective labour movements, were finding it difficult to implant themselves solidly and seemed condemned to lead a very marginal and precarious existence in many countries and for many years to come. But the Twenty-One Conditions remained the fundamental text of the Third International throughout its history and was used to serve the inquisitorial and arbitrary purposes of its centralized leadership in Moscow: the cost that was paid for this by the international working-class movement in the next twenty years and more is incalculable.

Leninism was a political style adapted - or at least intended to be adapted - to a particular political strategy - the political strategy that I have called 'insurrectionary politics'. Stalinism, which soon after Lenin's incapacitation and death came to dominate Russia and the international Communist movement, made a frightful caricature of the style, and made of the strategy what it willed, depending on what Stalin wanted to achieve at any particular time and in any particular place. In fact, Leninism as a coherent strategy of insurrectionary politics was never seriously pursued by the Third International, which means of course that it was never seriously pursued by its constituent Communist parties.

Nor was it pursued by Communist parties in the countries of advanced capitalism after the dissolution of the Comintern in 1943. Of course, the leadership of these parties remained wholly subservient to Stalin's command after its dissolution. But it is not really plausible to attribute these parties' renunciation of insurrectionary politics—for this is in effect what occurred—to Stalin and Stalinism. The renunciation was no doubt gre'atly encouraged by Moscow. But the reasons for it have much deeper roots than this, as is testified by the weakness of the opposition to this renunciation, inside the Communist parties and out. If it had not corresponded to very powerful and compelling tendencies in the countries concerned, the abandonment of the strategy of insurrectionary politics would have encountered much greater resistance in revolutionary movements, notwithstanding Moscow's prestige and pressure and repression. Much more than this was here at work, which condemned left opposition to Moscow to extremely marginal significance, and which is crucially important for the understanding of the life and character of the labour movements of advanced capitalism.

The essential starting-point is not only that the revolution which the Bolsheviks so confidently expected did not occur in advanced capitalist countries at the end of the First World War and that is was never even much of a realistic prospect; but that the old leaderships which the Third International and its constituent national units wanted to supplant and destroy remained m command of a large and in most cases the largest part by far of their labour movements, industrial as well as political. It is true that these leaderships had to make certain verbal and programmatic concessions to the radical temper that had been generated by the war and its aftermath. But this really amounted to very little in practical terms.

In fact, reformist leaderships, now confronted with a Bolshevik and Communist presence which they hated and feared as much as their bourgeois and conservative counterparts, became even more 'reformist' than they had been. For all practical purposes, the old parties of the now defunct Second International were, evenmore strongly after 1918 than before, parties of social reform—and parties of government as well. All in all, the leaderships of these parties withstood the Communist challenge with remarkable and significant success. Another way of putting this is to say that, in the sense in which the term has been used here,  the politics of Leninism, insurrectionary politics, failed in the countries of advanced capitalism. What was left in these countries from the upheaval of war and revolution was, on the one hand, parties of social reform whose leaderships had absolutely no thought of revolution of any kind, and who indeed saw themselves and their parties as bulwarks against it; and, on the other hand, Communist parties formally dedicated to revolution yet pursuing at the behest of the Comintern opportunistic and wayward policies, with more or less catastrophic consequences.

Nowhere was a socialist 'third force', distinct from the other two, able to achieve any substantial support. Indeed, this is putting it too strongly: it is more accurate to say that nowhere was a Marxist left, independent from both the parties of the old Second International and of those of the new Third International able to play more than a very marginal role in the life of their labour movements. The point is particularly relevant and significant in relation to the Trotskyist opposition which, from the second half of the twenties onwards, could plausibly lay claim to being the heir of Leninist insurrectionary politics. Trotskyism had an illustrious leader and small groups of dedicated and talented adherents in many capitalist countries. But these groups remained utterly isolated from their working-class movements and were never able to mount a serious challenge to their respective Communist parties.

Many different reasons have been advanced from within Marxism to account for the disappointments and defeats which soon came to mock the high hopes of 1917 and after. These reasons have to do with the unexpected capacity of capitalism to take the strain of economic dislocation and slump; the equally unexpected capacity of conservative forces to defend their regimes by ideological manipulation and, when necessary, by physical repression, with fascism and Nazism as the most extreme forms of that defence; and also the capacity of capitalism to respond to crisis and pressure with cautious and piecemeal reforms dressed up in a rhetoric of wholesale renewal—for instance the New Deal in the United States\index{America!the New Deal}. And Marxists also pointed to the role of social democratic leaders as, in effect, the defenders of capitalism and the agents of social stabilization—although Communists ceased to make as much of this in the thirties, when the Comintern adopted the Popular Front strategy, as they had done earlier.

Any such list of explanations why Leninist politics failed to make moie headway for this is the question—needs however to include one other factor of the greatest importance: this is the extremely strong attraction which legality, constitutionalism, electoralism, and representative institutions of the parliamentary type have had for the overwhelming majority of people in the working-class movements of capitalist societies.

Lenin himself, as it happens, was very concerned to stress the importance of this in his polemic against left Communists in 1920. In Western Europe and America\index{America}', he wrote, 'parliament has become most odious to the revolutionary vanguard of the working class'. But this, he went on to warn, was not true of the mass of the population of capitalist countries: . in Western Europe', he also wrote, 'the backward masses of the workers and —to an even greater degree—of the small peasants are much more imbued with bourgeois-democratic and parliamentary prejudices than they were in Russia\dots{}'\postnote{'Left-Wing' Communism — An Infantile Disorder, in SWL, pp. 549, 551.}

Even this, however, grossly underestimated the extent and tenacity of what Lenin called 'bourgeois\hyp{}democratic and parliamentary prejudices'; which is no doubt why he was led to make such absurd pronouncements as that 'in all civilised and advanced countries\dots{} civil war between the proletariat and the bourgeoisie is maturing and is imminent.'\postnote{Ibid., p. 548.}

Marxism and revolutionary politics had to adapt—or at least id adapt to the fact that Lenin was absolutely wrong. One form which that adaptation took was a new stress on 'socialism m one country', which soon became part of Stalinist dogmatics. In a wider perspective, Marxist politics continued to have not only one meaning but two. This perpetuated the situation which had existed when Leninism as insurrectionary politics had affirmed itself as an alternative strategy to the 'reformism' of the Second International: but it was now the Third International and its Communist parties which adopted the 'reformist' strategy.

What occurred in the early thirties, after the Comintern's catastrophic 'Third Period' and with the adoption of the Popular Front Pollcies, was the barely acknowledged and theoretically still unformulated abandonment by international Communism of insurrectionary politics. This, as I have argued earlier in discussing the meaning of 'reformism', did not involve the abandonment of a programme of wholesale transformation of capitalism, of class struggle on different fronts, of the possibility  that conservative violence would have to be met with violence, and least of all the abandonment of unconditional support for the Soviet Union. In fact, this unconditional support gave to Communist parties a stamp of particularity which helped to further the impression that they remained committed to insurrectionary politics long after this had ceased to be the case.\footnote{This unconditional support of the Soviet Union even thrust some Communist parties into illegality at certain times, for instance the French Communist Party which came under proscription at the beginning of the Second World War for its active opposition to the war.} Nor did Communist 'reformism' require the formal and explicit abandonment by Communist parties of their doctrinal commitment to Marxism, or for that matter to Marxism-Leninism.\footnote{Thus it is only in 1976 that the French Communist Party explicitly abandoned the concept of the dictatorship of the proletariat, which it had in fact long abandoned implicitly. For the meaning and significance of this abandonment, see below.}
 
But it did involve a progressively more definite acceptance of constitutionalism and electoralism as a possible and desirable 'road to socialism' for Communist parties, and the rejection of any notion of a 'seizure of power', at least in circumstances where legality and constitutionalism were so to speak available, or had come to be available—and even in many cases where it was not.

This strategy was developed and formulated over a fairly protracted period of time. It was not pursued with full explicitness in its earlier stages and there were cases where Stalin decided that it should not be pursued at all—thus in Eastern Europe, but under obviously exceptional conditions. In the liberated countries of Western Europe at the end of the Second World War, and already in Italy in 1943, it was pursued very faithfully.\footnote{In a very different situation and for different reasons, the strategy had already been applied by the Spanish Communist Party in the Spanish Civil War, under Russian direction.} The French and Italian Communist leaderships, which then acted only with the accord of the Russian leaders, opted decisively and overwhelmingly for constitutionalism and electoralism. They entered coalition governments in an essentially subsidiary capacity so far as actual policy-making was concerned, but played a critical role in the stabilization of their capitalist regimes at a time of great crisis for these regimes.

Whether a different strategy, of a more 'revolutionary' kind, would then have been realistic and capable of yielding better results has long been argued over and cannot obviously be settled conclusively. More important is the fact that the strategy that was then pursued, and which was the logical continuation of a process which the war had only partially suspended was continued and further developed in the following years, notwithstanding all rebuffs, vicissitudes, and disappointments.\footnote{The first Communist Party to give explicit articulation to this 'road to wTs firT TP ,Bntlsh' whose Programme, The British Road to Socialism, was first published m 1951, with the full approval of Stalin.} The French Communist Party was suddenly offered the opportunity to move in the direction of insurrectionary politics in May 1968 and rejected it with little if any hesitation. Here too, there has been endless argument as to whether a chance of revolution was then missed. But whichever view may be taken of this, there is at least no question that the P.C.F. very firmly acted on the conviction that any resort to insurrectionary politics would have been an utterly irresponsible form of 'ultra-left' adventurism which would have been crushed and would have set the working-class movement back for decades.

The spelling out of this 'reformist' strategy gathered quickening pace in the following years, particularly in Italy and then France, and is now thoroughly articulated: advance by way of electoral gains at all levels and notably at the national level; alliance with other left-wing or radically-inclined parties (though not usually with the 'ultra-left' parties and groupings) for the purpose of obtaining an electoral and parliamentary —and where appropriate a presidential—majority; the formation of a coalition left-wing government in which the Com- munists must play an important and possibly predominant role; and the accomplishment of a vast programme of social and economic reforms designed to begin the structural transformation of capitalist society and the eventual achievement of a socialist society. In the course of this process—and what is involved is a process, which is seen as stretching over a number of years and more-a plurality of parties would be maintained, and this would include anti-socialist parties; civic freedoms—of speech, association, movement, etc.-would not only be preserved but enhanced; representative institutions would not only continue to function but would be reformed so that they might function more democratically and responsively; and the executive power, owing its legitimacy to universal suffrage, would only seek to retain it on the same basis, by way of free and unfettered elections.

Another and major ingredient of this strategy concerns foreign policy, and particularly relations with the U.S.S.R. The stress has been increasingly on complete national independence and freedom from any kind of subservience to Soviet Communism. At most, Western Communist parties envisage that a government of which they were part, and even one of which they were the major part, would have no more than a friendly relation with the Soviet Union—but not to the detriment of friendly relations with other countries or even to the detriment of existing alliances such as NATO. A government in which Communists were involved might seek to end the system of alliances created at the time of the Cold War—NATO and the Warsaw Pact. But this is a very different matter from seeking a unilateral withdrawal from NATO. Whatever else 'proletarian internationalism may mean for these Communist parties nowadays, it means neither automatic support for the Soviet Union's external policies, nor abstention from criticism of its internal policies, notably for its derelictions in the realm of civic and democratic freedoms.

The fact that Communist parties in capitalist countries, with varying qualifications that do not affect the main drift, should have come to adopt so definite a 'reformist' strategy is obviously a matter of vast importance in general political terms as well as being important in relation to Marxist political theory; and this must be examined further.

In essence, two main questions are raised by the 'reformist' strategy; the first, and by now the less important one, is whether the achievement of executive power by electoral means is possible. The second, which is of a different order altogether, is what happens when such power has been achieved by Communists, or by Communists in a left coalition. Clearly, these are central questions in modern Marxist politics.

\section*{3}
It was at one time fairly common for Marxists to argue that, if it ever looked as if a left-wing party or parties pledged to radical anti-capitalist policies might win an electoral majority, the ruling class would not allow this to happen. Something like a pre-emptive strike , it was argued, would be launched, most probably from within the state system and by way of a military coup. The elections would not be held; and the constitutional regime which had brought about the electoral threat would be abrogated in favour of one form or other of right-wing dictatorship.

It was very reasonable for Marxists to treat this as one possibility, and not only in cases where universal suffrage appeared to pose a threat of fundamental change: the same might also happen where elections threatened to produce a government pledged to carry out reforms which, though not actually subversive of the existing social order, were judged dangerous by people in power. After all, this is precisely what happened in Greece in 1967, when the Colonels took over in anticipation of a general election which the Right feared it would lose and which would then have brought to power a mildly reforming government. The propensity of the Right to try and mount pre-emptive 'strikes' of this kind must vary from one country to another, and depend on many different factors: but it is reasonable (and prudent) to see it always and anywhere as one possibility—and the greater the danger, from the point of view of the Right, the greater the possibility.

It is clearly not reasonable, on the other hand, to treat the possibility as a certainty. For there are many circumstances which may well and in some cases almost certainly will render a pre-emptive coup impossible, or turn it into a gamble too desperate to be actively considered by serious people. Particularly in countries with strong electoral traditions, and with well- implanted labour movements, the chances are that, in more or less normal conditions, 'dangerous' elections will be held, and may bring to office left-wing governments in which Communists would have a substantial representation and possibly a preponderant one, and whose announced intention was to bring about the fundamental transformation of capitalist society.

It is at this point, however, that there occurs a crucial theoretical split in Marxism between 'reformism' on the one hand and its Leninist alternative on the other. The question at issue is the 'dictatorship of the proletariat' and its institutional meaning.

In the last chapter, I stressed the importance which Marx (and Engels) had attached to the notion that 'the working class cannot simply lay hold of the ready-made state machinery, and wield it for its own purposes'; and I also noted the reaffirmation by Lenin  of the meaning of that notion after its neglect by the Second International, in terms of the 'smashing' of the existing state power and its replacement by 'other institutions of a fundamentally different type'.\postnote{See above, pp. 138 ff.}

In The State and Revolution, Lenin appeared to mean by this phrase popular power expressed directly or through the Soviets or Workers and Soldiers' Deputies. Another and more specific example of this search for new institutional forms was provided by Rosa Luxemburg in the draft programme which she wrote for the new German Communist Party and which appeared as an article in Die Rote Fahne on 14 December 1918, under the title 'What Does the Spartacus Union Want?'
\begin{displayquote}
    The main feature of the socialist society is to be found in the fact that the great mass of the workers will cease to be a governed mass, but, on the contrary, will itself live the full political and economic life and direct that life in conscious and free self-determination.
\end{displayquote}

Therefore, she went on,
\begin{displayquote}
    the proletarian mass must substitute its own class organs—the workers' and soldiers' councils—for the inherited organs of capitalist class rule—the federal councils, municipal councils, parliaments—applying this principle from the highest authority in the state to the smallest community. The proletarian mass must fill all governmental positions, must control all functions, must test all requirements of the state on the touchstone of socialist aims and the interests of its own class;\postnote{H. Gruber, op. cit., p. 106.}
\end{displayquote}

and the programme also specified that workers' councils would be elected all over Germany, as well as soldiers' councils, and that these would in turn elect delegates to the Central Council of the workers' and soldiers' councils; and the Central Council would elect the Executive Council as the highest organ of legislative and executive power.\postnote{Ibid., p. 111.}

The economic and social points of this programme were of an equally far-reaching nature;\footnote{Thus the programme of economic demands included the expropriation of the land held by all large and medium-sized agricultural concerns and the 'establishment of socialist agricultural cooperatives under a uniform central administration all over the country'; the nationalization by the Republic of Councils of all banks, ore mines, coal mines, as well as all large industrial and commercial establishments; the confiscation of all property exceeding a certain limit; the take over of all public means of transport and communication; the election of administrative councils in all enterprises, 'to regulate the internal affairs of the enterprises in agreement with the workers' councils'; etc.} and the split between 'reformism' and what I have called its Leninist alternative undoubtedly includes; very considerable differences in the scale and extent of the immediate economic and social transformations which each strategy tends to propose. But the crucial theoretical difference consists in the acceptance on the one hand of the notion that the existing bourgeois state must be 'smashed' and replaced by an altogether ditferent type of state embodying and expressing the 'dictatorship of the proletariat'; and on the other hand, the more or less explicit rejection of any such notion.

It might seem as if the difference stemmed from a perspective of peaceful and constitutional transition as opposed to a violent one, the first type occurring on the basis of a left-wing victory at the polls and the second on the basis of a seizure of power made possible by a regime's defeat in war, or extreme economic crisis and dislocation, or political breakdown, or some combination of any of these possibilities. But this is not in fact where the opposition necessarily lies: a constitutional accession to power might be followed by a wholesale recasting of state institutions; and a seizure of power need not involve such a recasting at all. Indeed, in so far as the intention is to extend popular power, a peaceful transition might be more favourable to such a project than a violent one.\postnote{See below, pp. 188-9.}

But in any case, the essential theoretical difference is between a project which envisages the carrying through of a socialist ransformation by way of the main political institutions— no ably parliament— inherited from bourgeois democracy, even ough these might be to a greater or lesser extent reformed in more democratic directions; and a project which envisages the total transformation of the existing political institutions (i.e. the smashing of the state) as an integral and essential part of a socialist revolution. The classical Marxist text for this contraposition is Lenin's The Proletarian Revolution and the Rene- gade Kautsky which he wrote in 1918 in reply to Kautsky's The Dictatorship of the Proletariat (itself a critique of Lenin's The State and Revolution).

These two 'scenarios' or 'models'—'reformism' on the one hand and its Leninist alternative on the other-have constituted the two poles of a debate which has been at the centre of Marxist po ltics for the best part of this century; and the 'reformist' strategy to which Communist parties now firmly subscribe ensures that this debate will be sustained for a long while yet,  given the bitter opposition which that strategy evokes from various other tendencies of the Marxist left. It may well be, however, that the terms of the contraposition are mistaken, in so far as neither 'model' represents realistic perspectives and projections. I believe, on a number of different grounds, that this is indeed the case: whatever the real differences between the two strategies, they cannot be as stated, because the stated positions do not correspond to any possible situation that may be envisaged.

To begin with the Leninist strategy, it is by now clear (or should be) that the 'smashing' of the existing bourgeois state does not open the way for the achievement of the 'dictatorship of the proletariat'. A revolutionary situation produced by successful insurrectionary politics as a result of which the existing state is 'smashed' requires, if the revolution is to succeed and be defended and consolidated, a new articulation of power of a kind which cannot be provided by the 'dictatorship of the proletariat'—at least not in the meaning which Marx may be taken to have given to it, or which Lenin gave to it in The State and Revolution and even in The Proletarian Revolution and the Renegade Kautsky.

In that meaning, it is always popular power which is stressed, and it is its exercise by the proletariat and its allies which is glorified. That power is exercised almost directly, and only mediated by representative institutions strictly subordinated to the people. In The State and Revolution, Lenin thus writes that the state, under these conditions, is transformed 'into something which is no longer the state proper'.\postnote{V. I. Lenin, The State and Revolution, in SWL, p. 293.} It is necessary to suppress the bourgeoisie and crush their resistance: 'the organ of suppression, however, is here the majority of the population, and not a minority, as was always the case under slavery, serfdom and wage slavery.\postnote{Ibid., p. 293.} He then goes on:
\begin{displayquote}
    and since the majority of the people itself suppresses its oppressors, a 'special force' for suppression is no longer necessary! In this sense, the state begins to wither away. Instead of the special institutions of a privileged minority (privileged officialdom, the chiefs of the standing army), the majority itself can directly fulfil all these functions, and the more the functions of state power are performed by the people as a whole, the less need there is for the existence of this power.\postnote{Ibid., p. 293.}
\end{displayquote}

True, Lenin says, 'we are not Utopians, we do not “dream” of dispensing at once with all administration, with all subordination.' But this subordination 'must be to the armed vanguard of all the exploited and working people, i.e. to the proletariat'.\postnote{Ibid., p. 298.}

A similar trend of thought is evident in The Proletarian Revolution and the Renegade Kautsky. Thus, 'the revolutionary dictatorship of the proletariat is rule won and maintained by the use of violence by the proletariat against the bourgeoisie, rule that is unrestricted by any laws';\postnote{V. I. Lenin, The Proletarian Revolution and the Renegade Kautsky, in CWL, vol. 28 (1965), p. 236.} and Lenin insists that 'proletarian democracy is a million times more democratic than any bourgeois democracy' and 'Soviet power is a million times more democratic than the most democratic bourgeois republic' because it was based on the Soviets: 'The Soviets are the direct organisation of the working and exploited people themselves, wnich helps them to organise and administer their own state in every possible way.'\postnote{Ibid., pp. 246-7.}

There is, however, an important difference between the first pamphlet and the second. By the time Lenin wrote the second one, he was leading a state that had not been transformed 'into something which is no longer the state proper', and would not be, and could not be.

In its proper Marxist meaning, the notion of the 'dictatorship of the proletariat' disposes much too easily, and therefore does not dispose at all, of the inevitable tension that exists between the requirement of direction on the one hand, and of democracy on the other, particularly in a revolutionary situation. As I noted in the last chapter, Lenin met the problem by affirming that the dictatorship of the party (and therefore of the state which the party controlled) was the dictatorship of the proletariat. But it did not take him long to realize that he had only redefined the term rather than fulfilled its original promise. In fact, it is scarcely too much to say that a time of revolution of the Leninist kind is precisely when the 'dictatorship of the proletariat' is least possible—because such a time requires the re-creation of a new and strong state, a 'state proper', on the ruins of the state which the revolution has 'smashed'. But this is not the 'dictatorship of the proletariat' of Marx, and of Lenin in The State and Revolution.

Nor can it be contended that the pyramid of councils, of the kind which Rosa Luxemburg projected in 1918, resolves the question of direction and democracy. For the structure which she proposed was, very reasonably, to be capped by an Executive Council as the highest organ of legislative and executive  power';\postnote{Gruber, op. cit., p. 111.} and it is clear that, whatever her intentions, this organ would have been, and would have had to be, an extremely strong state, possibly representing the dictatorship of the proletariat', but not amounting to it.

Where power has been seized, revolutionaries have to create a strong state in place of the old if their revolution is to survive and begin to redeem its promise and purpose. This is bound to be an arduous task, particularly because the material circumstances in which it has to be undertaken are likely to be unfavourable and further aggravated by the hostility and opposition of the new regime's internal and external enemies. Inevitably some of its own supporters, and possibly many, will falter and turn away when the exaltation of the first phase wears off as it confronts the mundane and difficult requirements of the second.\footnote{Gramsci put the point about these requirements rather well: 'What is needed for the revolution', he wrote, 'are men of sober mind, men who don't cause an absence of bread in the bakeries, who make trains run, who provide the factories with raw materials and know how to turn the produce of the country into industrial produce, who ensure the safety and freedom of the people against the attacks of criminals, who enable the network of collective services to function and who do not reduce the people to despair and to a horrible carnage. Verbal enthusiasm and reckless phraseology make one laugh (or cry) when a single one of these problems has to be resolved even in a village of a hundred inhabitants' (A. Gramsci, Ordine Nuovo, 1919-1920 (Turin, 1953), pp. 377—8) in N. Badaloni, Gramsci et le Probleme de la Revolution' in Dialectiques, nos. 4—5, March 1974, p. 136. The quotation is from an article written in June 1919.} The new regime may retain a very wide measure of popular support and find it possible to rely on continued popular involvement. It will most probably go under quite soon if it cannot. But the tension remains between state direction and popular power; and that tension cannot be resolved by invocations and slogans. Gramsci said in 1919 that
\begin{displayquote}
    The formula 'dictatorship of the proletariat' must cease to be a formula, an occasion to parade revolutionary rhetoric. He who wills the end, must will the means. The dictatorship of the proletariat is the installation of a new State, typically proletarian, into which flow the institutional experiences of the oppressed class, in which the social life of the worker and peasant class becomes a system, universal and strongly organised.\postnote{L'Ordine Nuovo, 21 June 1919, inNewEdinburgh Review. Gramsci. II, p. 54.}
\end{displayquote}

This at least acknowledges the central place which a 'state proper must play in the new system. But in so doing, it only proposes an agenda, not a blueprint.

The 'reformist' strategy which is counterposed to Leninism and the 'dictatorship of the proletariat' by Communist parties (though this is not of course how they themselves put it), raises a different series of questions.

Let it be assumed that a coalition of left-wing parties, in which the Communist Party has an important or preponderant place, wins a general election on a programme of pronounced anticapitalist policies, and 'is allowed' to come to office. What happens then?

Left critics of the 'reformist' strategy have usually argued that the answer to this was—not much. In other words, 'reformist' leaders would not seek to implement the programme on which they had been elected. No doubt there would be a good many changes in the personnel of the existing state system—not least because an hitherto excluded political army would expect to be provided with rewards, at all levels of the administrative apparatus. Some institutional and administrative reforms might be attempted. Some social reforms would be proposed and possibly implemented; and even some measures of state ownership would be introduced. But this is as much as could be expected from leaders, so the argument goes, who are by now well integrated in their bourgeois political systems, and who would act as agents of stabilization of the existing social order rather than as architects of a new one. In their role as agents of stabilization, it is further said, they would need (and would be willing) to contain and repress the militancy of the working class, which their coming to power would itself have enhanced. Thus the self-proclaimed Party of Revolution would turn into the shamefaced Party of Order—and perhaps not so shamefaced at that; and in due course, it would be thrust out of office, with the old order more or less intact.

This is a very possible 'scenario'. It might need amendment here and there; but it is certainly not unreasonable to think that some such outcome could be produced by the coming to office of a left-wing government with Communist participation. After all, Communist participation in governments in Italy and France and other countries at the end of the war and immediately after did produce precisely such an outcome. Admittedly, the governments in question were not left-wing ones but bourgeois coalitions in which the Communists accepted a subsidiary role.

But while circumstances would now be different, the interven-  ing years have also witnessed a certain 'social-democratization' of Communist parties, which is likely to be enhanced by the pressures of office, and which is unlikely to be compensated for by pressures from the left inside and outside these parties. This makes it not at all absurd for intelligent conservatives (as well as critics on the left) to believe that a left-wing government with substantial Communist participation—or even Communist- led—might not represent nearly so great a threat to capitalism and the existing social order as less intelligent conservatives believe, and that it might even turn out to be a necessary price to pay for the maintenance of the system.

But let it be supposed, on the other hand, that a left\hyp{}wing and Communist\hyp{}led government, backed by an electoral and parliamentary majority, did decide to carry through far\hyp{}reaching anti\hyp{}capitalist measures, including fairly drastic measures of nationalization, imperative economic co\hyp{}ordination and planning, major advances in labour legislation, social welfare, education, housing, transport and the environment, all favouring the 'lower income groups' and with fiscal policies directed against the rich. This is presumably what should be expected from a government whose purpose was claimed to be the radical transformation of capitalist society in socialist directions. It might not be able to do everything at once; but it would at least be expected to make a convincing beginning. What then?

Perhaps the first thing to note about this 'scenario' is that there are very few precedents indeed for it. In fact, the only precedent which comes reasonably close to it is that of Salvador Allende\index{Allende} in Chile. This is not of course to say that what happened in Chile under Allende is conclusive, though it is significant.\postnote{See R. Miliband, 'The Coup in Chile', The Socialist Register 1973 (London, 1973). 1 use some of the ideas and concepts put forward in that article in what follows here.
} It is rather to point to the fact that, for all the controversies which have surrounded 'reformism', its strategy has hardly ever been put to the test: such 'reformist' governments as there have been have taken the first of the two options, and done nothing much—or even nothing at all—to cause concern to the powers-that-be.

The next thing to be said about the attempt to carry through the kind of measures listed above is surely that it would be bound to arouse the fiercest enmity from conservative forces defeated at the polls but obviously very far from having lost all their formidable class power.

This class power endures both inside the state system and in society at large. Of course, it must be assumed, as part of the 'scenario', and as partially demonstrated by the electoral victory of the forces of the left, that these forces of the left have deeply permeated society, that they are well implanted in most spheres of life, and that they have many friends in the state system. But to take the latter first, it is obviously realistic to expect that by far the larger part of the state personnel at the higher levels, and at least a very large number in the lower ones as well, are much more likely to be ideologically, politically, and emotionally on the side of the conservative forces than of the government. In many cases, they will only be suspicious and hesitant. But in many others, they will be firmly opposed to programmes and policies which they believe to be utterly detrimental to the 'national interest'. Many civil servants throughout the bureaucracy, many regional and local administrators and officials, and most members of the judiciary at all levels, together with the police and the officer class in general, may—and indeed must in prudence—be taken to be in varying degrees opposed to the government. The commitment need not be thought to be irrevocable: but it must be assumed to exist at the start.

No doubt new people, attuned to the government's purposes, will be brought in; and new administrative organs, staffed by such people, may be set up. But it is just as well not to underestimate the degree of administrative dislocation and even chaos that may be involved in the process, and which people of conservative disposition in the state system have no reason to try and remedy—rather the reverse. This means that a battle—in fact the class struggle—will also be waged inside the state system at all levels; and it would, incidentally, have to be waged, even if the state had previously been 'smashed'. But how that battle goes depends to a large extent on what happens outside the state system as well as inside it. Those involved on the conservative side will be deeply influenced in their thinking and behaviour by the manner in which the forces of the left, beginning with the government, are waging the struggle—by their determination, intelligence, and good sense; but also, and crucially, by their capacity and will to rely on popular support and initiative, of which more in a moment.

The class power of the conservative forces in capitalist society at large assumes many different forms—the control of strategic means of industrial, commercial, and financial activity, and of vast resources; the control of most of the press, and many other  means of political communication and of ideology in general; large, well-implanted political parties, associations, pressure groups and organizations of every sort, many of which pride themselves on their entirely 'non-political' character, meaning that they are conservative without being affiliated to any particular party; and to a greater or lesser degree, depending on the particular country but nowhere negligible, the churches and their satellite organizations.

Nor is it only a matter of organized, collective power; it is also one of influence and activity on the part of individuals, each making a contribution in his and her own sphere to the strengthening of the conservative forces and the advancement of conservative purposes. It must be reckoned that large numbers of men and women belonging to the middle and upper classes, and occupying positions of relative influence and responsibility in their communities, will be willing to help, in whatever way they can, in the task of saving the country, defending freedom, national independence, their children's future, or whatever.

It may be assumed that a process of acute polarization will have occurred in the weeKs and months preceding the accession of the new government to office, and notably during the electoral campaign. But the struggle enters an entirely new phase once the government is in office. There may be a short respite, because of the sense of disappointment and demoralization which besets the defeated side. But this will not last long, and the conservative forces will soon be reorganizing themselves for a task which appears every day more pressing, namely the destabilization of a hated government, and its ultimate undoing. In this task, they will have the precious assistance of powerful capitalist governments and international capitalist interests. As I have noted in a previous chapter, this dimension of the class struggle is of quite vital importance. Indeed, so much is this the case that any socialist 'experiment' of the kind envisaged here would undoubtedly depend very greatly on international solidarity and support, capable of neutralizing or at least of attenuating the impact of the efforts at destabilization that would be undertaken by internal and external conservative forces. Much of the responsibility for countering these latter forces would depend on working-class movements in advanced capitalist countries. In many instances, this would involve a difficult struggle against social-democratic leaders who must, on any realistic view, be reckoned as forming part, and a very important part, of the conservative forces of advanced capitalist countries. The notion that a strongly reforming government in one such European country would mainly face the opposition of the United States is now seriously out of date: there are other countries—Germany and Britain, for instance—with social- democratic governments at their head which could be expected to help in trying to bring the offending government to heel. The responsibility of the left in the countries concerned would clearly be to make their task more difficult, and preferably to make it impossible.

Every reforming step that the government takes will reinforce the determination of its opponents to see it defeated. The opposition will naturally use all the constitutional devices and institutional opportunities which are available to it, for instance in parliament. But this is only one aspect of the struggle and by no means the most significant.

Reference has already been made to the struggle which will be waged in the state system; and what happens here may well be of the greatest importance. There is much that a sullen and hostile bureaucracy can do to impede and discredit a radical and reforming government.

But the struggle will also be waged in every part of civil society—in factories and power stations, dockyards and warehouses, shops and offices, barracks, schools and universities; as well as the press, radio, and television; and also the streets. It will assume an infinite variety of forms, because all forms of social life become 'politicized' in circumstances of great social stress and crisis; or to put it more accurately, their 'political' character, instead of being as hitherto blurred by the general acceptance of prevailing ideas and values, becomes visible and obtrusive by virtue of general contestation, and that 'political' character is further sharpened and extended thereby.

It must be taken, by definition, that a 'reformist' government will try to keep this intensified class struggle within more or less constitutional bounds, and will hope to prevent it from spilling over into class war. It must expect some Fascist-type groups to engage in sporadic acts of violence; but it will hope that this can be confined and contained, and that the conservative forces in general will not encourage, support, let alone initiate and  foment unconstitutional action designed to overthrow the gov- ernment or m one way or another get rid of it.

But what the conservative forces in general decide to do or not to do must depend in a substantial degree on the government's own attitudes and actions. The fact that it has reached office by way of an electoral victory and that it has constitutional legiti- macy on its side gives it a psychological and political advantage which it would be absurd to under-estimate. But it would be equally absurd for such a government to expect that this will necessarily keep its opponents on the path of constitutional rectitude. Some elements among them, and possibly many, who had up to then been constitutionally-minded-since occasion had not arisen for them to be otherwise—will come to find it increasingly difficult if not impossible to keep on that constitutional path, and will persuade themselves that their patriotic uty requires them to stray from it and to encourage others to stray from it.

But the question is not simply one of constitutionalism and legality on the one hand and unconstitutionality, violence, military coups, and civil war on the other. What is involved, at east in the earlier stages and possibly over a prolonged period, is not outright violence and its encouragement, but economic! administrative, and professional forms of disruption and dislocation which may not be illegal at all and which are certainly not violent. What the government here faces is the pursuit by the conservative forces, with all the means at their command of many different forms of class struggle, this being sufficiently militant in its intention, character, and consequence to warrant its designation as class war rather than class struggle, yet falling well short of outright resort to arms and civil war.

In this situation, how the government responds is crucial. But it is also the sort of situation which a 'reformist' strategy finds it most difficult to handle. 'Reformist' leaderships know perfectly well that Marx was right when he said that universal suffrage may give one the right to govern but does not give one the power to govern. However, there is a great difference between knowing this and knowing how to act upon the knowledge and being willing to act upon it.

The government will be strongly pressed to 'be reasonable', to seek compromise and conciliation, to put its more 'extreme' proposals into cold storage until a 'more favourable' time, to consolidate and pause; and there will be those in its midst who will be very tempted to yield to such advice, and who will readily find an apt quotation from Lenin's 'Left-Wing' Communism—An Infantile Disorder to back up their advocacy of compromise, or an apt precedent from Bolshevik history to the same effect.

Nor is it necessarily the case that such advocacy is always mistaken - though the chances are that if it is accepted as a general strategy, it will only be interpreted by the conservative forces as a sign of weakness and indecision, and encourage them to press on with even greater determination their endeavours to destabilize the government; while the government's supporters will be discouraged and demoralized, confused and divided; and the large mass of waverers will be rendered more receptive to the appeals of the opposition.

If however the government does decide to push forward with its programme, it must also as a sine qua non strengthen itself in order to make advance—and survival—possible. Advance and defence are in this instance one and the same thing.

In effect, the government has only one major resource, namely its popular support. But this support, expressed at the polls, has to be sustained through extremely difficult times, and it has to be mobilized. The parties supporting the government will no doubt do this, or try to do it, in regard to their own members; and other working-class organizations, such as trade unions, will play a part. What is required, however, is something very much larger than can be provided by such organizations, namely a flexible and complex network of organs of popular participation operating throughout civil society and intended not to replace the state but to complement it. This is an adaptation of the concept of 'dual power', in that the organs of popular participation do not challenge the government but act as a defensive-offensive and generally supportive element in what is a semi-revolutionary and exceedingly fraught state of affairs.

In this perspective, a 'reformist' strategy, if it is taken seriously and pursued to its necessary conclusion, must lead to a vast extension of democratic participation in all areas of civic life—amounting to a very considerable transformation of the character of the state and of existing bourgeois democratic forms. If this be so, it turns out that the 'reformist' strategy ultimately involves the acknowledgement of the truth of the proposition of Marx and Engels that 'the working class cannot simply lay hold of the ready\hyp{}made state machinery, and wield it for its own purpose.'

The truth of that proposition does not, however, confer validity on the proposition which Marx and Engels and Lenin linked with it, namely that the existing state must be 'smashed' in order to replace it by the 'dictatorship of the proletariat'. As I have suggested earlier, the link that they established between the 'smashing' of the existing state and what they conceived to be the 'dictatorship of the proletariat' is an illusory one. What follows the 'smashing' of the existing state is the coming into being of another 'state proper', simply because a 'state proper' is an absolutely imperative necessity in organizing the process of transition from a capitalist society to a socialist one.

That process of transition both includes and requires radical changes in the structures, modes of operation, and personnel of the existing state, as well as the creation of a network of organs of popular participation amounting to 'dual power'. The 'reformist' strategy, at least in this 'strong' version of it, may produce a combination of direction and democracy sufficiently effective to keep the conservative forces in check and to provide the conditions under which the process of transition may proceed.

There are many regimes in which no such possibility exists at all; and where radical social change must ultimately depend on the force of arms. Bourgeois democratic regimes, on the other hand, may conceivably offer this possibility, by way of a strategy which eschews resort to the suppression of all opposition and the stifling of all civic freedoms. Such a strategy is full of uncertainties and pitfalls, of dangers and dilemmas; and it may in the end turn out to be unworkable. But it is just as well to have a sober appreciation of the nature of the alternative and not to allow slogans to take over. Regimes which do, either by necessity or by choice, depend on the suppression of all opposition and the stifling of all civic freedoms must be taken to represent a disastrous regression, in political terms, from bourgeois democracy, whatever the economic and social achievements of which they may be capable. Bourgeois democracy is crippled by its class limitations, and under constant threat of further and drastic impairment by conservative forces, never more so than in an epoch of permanent and severe crisis. But the civic freedoms which, however inadequately and precariously, form part of bourgeois democracy are the product of centuries of unremitting popular struggles. The task of Marxist politics is to defend these freedoms; and to make possible their extension and enlargement by the removal of their class boundaries.

\printpostnotes